\documentclass[aspectratio=169]{beamer}

\usetheme{Boadilla}
\setbeamertemplate{navigation symbols}{}

\usepackage[T1]{fontenc}
\usepackage[utf8]{inputenc}
\usepackage{lmodern}
\usepackage{booktabs}
\usepackage{array}
\usepackage{xcolor}

\definecolor{hmblue}{RGB}{22,74,124}
\definecolor{hmlightblue}{RGB}{233,241,249}
\definecolor{hmgold}{RGB}{173,122,34}
\definecolor{hmgray}{RGB}{70,70,70}
\definecolor{hmsoft}{RGB}{246,247,249}
\setbeamercolor{structure}{fg=hmblue}
\setbeamercolor{frametitle}{fg=hmblue}
\setbeamercolor{title}{fg=hmblue}
\setbeamercolor{normal text}{fg=black,bg=hmsoft}
\setbeamercolor{block title}{fg=white,bg=hmblue}
\setbeamercolor{block body}{fg=black,bg=hmlightblue}
\setbeamercolor{itemize item}{fg=hmblue}
\setbeamercolor{enumerate item}{fg=hmblue}
\setbeamerfont{title}{series=\bfseries,size=\Large}
\setbeamerfont{frametitle}{series=\bfseries,size=\large}
\setbeamerfont{block title}{series=\bfseries}

\setbeamertemplate{frametitle}{
  \vspace{0.2em}
  \insertframetitle
  \par
  \vspace{0.2em}
  {\color{hmgold}\rule{\linewidth}{0.8pt}}
}

\title{Hot Memory}
\subtitle{GPU Profiling and a New Outsourcing Model}
\author{William Dawson}
\date{\today}

\begin{document}

\begin{frame}
  \titlepage
\end{frame}

\begin{frame}{Research Question and Delivery Goal}
  \textbf{Research question}
  \begin{itemize}
    \item I want to port an MPI code to a low-memory GPU.
    \item I am afraid the code will run out of memory.
    \item The real question is not total allocation, but:
    \item \hspace{1em} Which kernels touch how much \textbf{hot memory} while they run?
    \item \hspace{1em} Which kernels are compute-heavy, and which are memory-traffic-heavy?
  \end{itemize}

  \vspace{0.5em}
  \textbf{Outsourcing goal}
  \begin{itemize}
    \item We want to outsource not just a tool, but a way of working.
    \item The deliverable is a container with the profiler, the dependencies,
          and the agent instructions needed to use it.
    \item The user should be able to drop into \texttt{claude} and continue
          the analysis collaboratively after delivery.
  \end{itemize}
\end{frame}

\begin{frame}{Methodology: How We Grab Hot Memory}
  \begin{block}{Definition}
    The hot memory of a kernel is the set of unique memory pages touched
    while that kernel executes.
  \end{block}

  \vspace{0.5em}
  What Linux is tracking for us:
  \begin{itemize}
    \item each virtual-memory page has a \texttt{Referenced} state
    \item if the process touches a page, the kernel marks it as referenced
    \item \texttt{/proc/self/smaps} exposes those per-mapping reference counts
  \end{itemize}

  \vspace{0.5em}
  What our instrumentation does:
  \begin{enumerate}
    \item \texttt{WSS\_BEGIN()} writes \texttt{1} to \texttt{/proc/self/clear\_refs}
    \item Linux clears the current referenced state on the process pages
    \item the kernel runs normally
    \item pages touched during that interval become referenced again
    \item \texttt{WSS\_END()} reads \texttt{/proc/self/smaps} and sums the
          new \texttt{Referenced:} totals
  \end{enumerate}

  \vspace{0.5em}
  So the result is:
  \begin{itemize}
    \item the unique 4 KB pages touched during one kernel
    \item not the whole process allocation
    \item exactly the working-set quantity we need for GPU residency decisions
  \end{itemize}
\end{frame}

\begin{frame}{Example: One Kernel Broken Down}
  \begin{block}{Example situation}
    The application allocates much more memory than any one kernel actually needs.
  \end{block}

  \vspace{0.5em}
  \begin{columns}[T,totalwidth=\textwidth]
    \begin{column}{0.47\textwidth}
      \textbf{What we want to know}
      \begin{itemize}
        \item Peak RSS might be 20 GB.
        \item This kernel may only touch 2 GB while it runs.
        \item That 2 GB is the footprint that must be near the GPU.
      \end{itemize}
    \end{column}
    \begin{column}{0.47\textwidth}
      \textbf{How to read it}
      \begin{itemize}
        \item If the kernel is memory-bound and 2 GB does not fit, that is bad.
        \item If it is compute-bound, some transfer cost may be hidden.
        \item If 2 GB fits, colder data may move outside the kernel.
      \end{itemize}
    \end{column}
  \end{columns}

  \vspace{0.5em}
  So each kernel gets three signals:
  \begin{itemize}
    \item \textbf{Hot MB}: what must fit
    \item \textbf{GFLOP}: arithmetic work
    \item \textbf{Memory accesses}: traffic / reuse pressure
  \end{itemize}
\end{frame}

\begin{frame}{Workflow Embedded in the Skill}
  The profiling skill is not just documentation. It encodes a way of working
  with the user.

  \vspace{0.5em}
  \begin{columns}[T,totalwidth=\textwidth]
    \begin{column}{0.5\textwidth}
      \textbf{Workflow}
      \begin{enumerate}
        \item confirm baseline build and run
        \item run a capability check
        \item measure peak RSS
        \item run perf hotspot discovery
        \item propose instrumentation points
        \item rebuild with profiling
        \item run the profiled code and interpret the result
      \end{enumerate}
    \end{column}
    \begin{column}{0.46\textwidth}
      \textbf{Interaction model}
      \begin{itemize}
        \item explain the next step
        \item explain why it matters
        \item say what remains
        \item ask permission before major actions
        \item use the project's build system if it exists
      \end{itemize}
    \end{column}
  \end{columns}

  \vspace{0.5em}
  This is the outsourcing idea in practice:
  \begin{itemize}
    \item the methodology is packaged
    \item the toolchain is packaged
    \item the collaboration protocol is packaged
  \end{itemize}
\end{frame}

\begin{frame}{Hardest and Ongoing Challenges}
  \begin{columns}[T,totalwidth=\textwidth]
    \begin{column}{0.31\textwidth}
      \textbf{Profiling inside a container}
      \begin{itemize}
        \item R-CCS cloud jobs run inside containers themselves
        \item \texttt{/proc/self/clear\_refs} and \texttt{perf\_event\_open}
              may be blocked by the outer container policy
        \item Need \texttt{--privileged} or capability grants that the
              cluster may not expose
      \end{itemize}
    \end{column}
    \begin{column}{0.31\textwidth}
      \textbf{Keeping the agent honest}
      \begin{itemize}
        \item Agent ignores warnings in output and moves on
        \item Trusts documentation over observed results
        \item Tendency to report 0 GFLOP as the answer rather than diagnose
        \item Skill must be written defensively: force reconciliation,
              not trust
      \end{itemize}
    \end{column}
    \begin{column}{0.31\textwidth}
      \textbf{The right human experience}
      \begin{itemize}
        \item The goal is not automation
        \item The goal is the optimal copilot: the user stays in control,
              the agent removes friction
        \item Interaction protocol matters as much as tool correctness
        \item Still learning what that balance looks like in practice
      \end{itemize}
    \end{column}
  \end{columns}
\end{frame}

\begin{frame}{Project}
  \begin{block}{GitHub}
    \texttt{https://github.com/william-dawson/hot-memory}
  \end{block}

  \vspace{0.5em}
  Repository contents:
  \begin{itemize}
    \item Singularity / Apptainer container definition
    \item WSS profiler library and MCP server
    \item profiler skill and code-skill template
    \item worked examples, including the built-in benchmark
  \end{itemize}
\end{frame}

\end{document}
